% ============================================================
% The Collision Spectrum and the L-Function Landscape
% ============================================================

\documentclass[12pt]{amsart}

\usepackage{amsmath,amssymb,amsthm}
\usepackage{booktabs}
\usepackage{microtype}
\usepackage{url}

\newtheorem{theorem}{Theorem}
\newtheorem{lemma}[theorem]{Lemma}
\newtheorem{proposition}[theorem]{Proposition}
\newtheorem{corollary}[theorem]{Corollary}
\theoremstyle{definition}
\newtheorem{definition}[theorem]{Definition}
\newtheorem{observation}[theorem]{Observation}
\theoremstyle{remark}
\newtheorem{remark}[theorem]{Remark}

\title{The Collision Spectrum\\and the $L$-Function Landscape}

\author{Alexander S.\ Petty}
\email{alexander.petty@gmail.com}
\date{November 2025 (revised April 2026)}
\subjclass[2020]{11A63, 11N05, 11M06}

\begin{document}
\begin{abstract}
The collision transform decomposes the centered
collision deviation $S^{\circ}$ over odd Dirichlet
characters modulo $m = b^{\ell+1}$. This paper proves
that at lag~$1$ (so $m=b^2$), for prime $b$ and
primitive odd $\chi$:
\[
\hat{S}^{\circ}(\chi)
= -\frac{B_{1,\overline{\chi}} \cdot
\overline{S_G(\chi)}}{\phi(m)},
\]
where $B_{1,\overline{\chi}}$ is the generalized
Bernoulli number and $S_G(\chi)$ is the diagonal
character sum. In magnitude:
$|\hat{S}^{\circ}| = |B_1| \cdot |S_G| / \phi(m)$.

Combined with the Bernoulli--$L$-value formula
\[
|B_1| = \frac{\sqrt{m}}{\pi}\,|L(1,\chi)|
\]
and the $S_G$ factorization theorem (proved in
a companion paper~\cite{paperC}), this gives
\[
|\hat{S}^{\circ}(\chi)|
= \frac{\sqrt{m}}{\pi\,\phi(m)}\,
|L(1, \chi)| \cdot |S_G(\chi)|.
\]
The collision invariant's Fourier spectrum encodes
$L$-function special values. Computation shows
$|S_G|$ itself correlates with $|L(1)|$, giving
$|\hat{S}^{\circ}| \sim |L(1)|^{\alpha}$ with
$\alpha = 2.00$ at base~$5$ (exact) and
$\alpha \approx 1.80$ across the tested prime bases.
\end{abstract}
\maketitle

% ============================================================
\section{Introduction}

The collision transform~\cite{paperB} decomposes the
centered collision deviation $S^{\circ}$ over Dirichlet
characters modulo $m = b^{\ell+1}$. The antisymmetry
theorem restricts the decomposition to odd characters.
The preceding papers recorded computationally that the collision
coefficients $|\hat{S}^{\circ}(\chi)|$ correlate with
$|L(s, \chi)|$ at Pearson $r > 0.87$ across all
prime bases tested~\cite{paper19,paper20}.

This paper proves the structural identity underlying
the correlation. The proof uses three ingredients: the
slice formula from~\cite{paperA}, the classical
Bernoulli identity for character sums over fractional
parts, and the vanishing of coset sums for primitive
characters.

% ============================================================
\section{The Decomposition Theorem}

\begin{theorem}[Decomposition]\label{thm:decompose}
Let $b$ be prime, $m = b^2$, and $\chi$ a primitive
odd character modulo~$m$. Then
\[
\hat{S}^{\circ}(\chi)
= -\frac{B_{1,\overline{\chi}} \cdot
\overline{S_G(\chi)}}{\phi(m)},
\]
where
$B_{1,\overline{\chi}} = (1/m) \sum_{a} a\,
\overline{\chi}(a)$ is the generalized first Bernoulli
number and
$G = \{r(b{+}1) : 0 \le r \le b{-}1\}$ is the diagonal
set, and
$S_G(\chi) = \sum_{n \in G}
[\overline{\chi}(n{+}1) - \overline{\chi}(n)]$
is the diagonal character sum.
\end{theorem}

\begin{proof}
Write $\phi = \phi(m)$. The Fourier coefficient is
\[
\phi\, \hat{S}^{\circ}(\chi)
= \sum_a S^{\circ}(a)\, \overline{\chi}(a)
= \sum_a \bigl[S(a) - \overline{S}_{R(a)}\bigr]\,
\overline{\chi}(a).
\]

\emph{Step 1: the centering term vanishes.}
The class mean $\overline{S}_{R}$ is constant on
each spectral class $R = (a{-}1) \bmod b$. The sum
over each class is
$\sum_{a \equiv k \bmod b} \overline{\chi}(a)
= \overline{\chi}(k) \sum_{j=0}^{b-1}
\overline{\chi}(1 + jb)$.
For primitive $\chi$ modulo $b^2$ (conductor $b^2$,
not $b$), the restriction of $\chi$ to the subgroup
$\{1 + jb : j = 0, \ldots, b{-}1\}$ is non-trivial.
By character orthogonality over this subgroup, the
sum vanishes. Therefore
$\sum_a \overline{S}_{R(a)}\, \overline{\chi}(a) = 0$.

\emph{Step 2: expand $S(a)$ via the slice formula.}
By~\cite{paperA},
$S(a) = -1 - \lfloor a/b \rfloor
+ \sum_{n \in G} d_n(a)$
where $d_n(a) = \lfloor(n{+}1)a/m\rfloor
- \lfloor na/m \rfloor$.

The constant term: $\sum_a (-1)\, \overline{\chi}(a)
= 0$ (non-trivial $\chi$).

The floor term: $\lfloor a/b \rfloor = a/b - \{a/b\}$.
Since $\{a/b\} = \{ba/m\}$ and the Bernoulli identity
(Lemma~\ref{lem:bernoulli} below) gives
$\sum_a \{ba/m\}\, \overline{\chi}(a)
= \chi(b)\, B_{1,\overline{\chi}} = 0$ (because
$\gcd(b, m) \ne 1$, so $\chi(b) = 0$):
\[
-\sum_a \lfloor a/b \rfloor\, \overline{\chi}(a)
= -(1/b) \sum_a a\, \overline{\chi}(a) + 0
= -b\, B_{1,\overline{\chi}}.
\]

\emph{Step 3: the diagonal terms.}
The endpoint slices are handled separately. For
$n = 0$, one has $d_0(a)=0$, so the contribution
vanishes. For $n = m{-}1$, one has
$d_{m-1}(a)=1$ for every unit $a$, so the sum against
the non-trivial character $\overline{\chi}(a)$ also
vanishes.

For the interior slices, Lemma~\ref{lem:bernoulli}
gives
$\sum_a \lfloor na/m \rfloor\, \overline{\chi}(a)
= [n - \chi(n)]\, B_{1,\overline{\chi}}$
for $\gcd(n, m) = 1$. Therefore
\begin{align*}
\sum_a d_n(a)\, \overline{\chi}(a)
&= [(n{+}1) - \chi(n{+}1)]\, B_1
- [n - \chi(n)]\, B_1 \\
&= [1 + \chi(n) - \chi(n{+}1)]\, B_1.
\end{align*}

Summing over $G$ ($|G| = b$ slices):
\[
\sum_{n \in G} \sum_a d_n(a)\, \overline{\chi}(a)
= B_1 \left[b - \sum_{n \in G}
(\chi(n{+}1) - \chi(n))\right]
= B_1 \bigl[b - \overline{S_G(\chi)}\bigr].
\]

\emph{Step 4: combine.}
\[
\phi\, \hat{S}^{\circ}(\chi)
= 0 + (-b\, B_1) + B_1\bigl[b -
\overline{S_G(\chi)}\bigr] - 0
= -B_1 \cdot \overline{S_G(\chi)}.
\]
Dividing by $\phi$ gives the result.
\end{proof}

\begin{lemma}[Bernoulli identity]\label{lem:bernoulli}\leavevmode\newline
For primitive $\chi$ modulo $m$ and $\gcd(n,m)=1$:
\[
\sum_{\substack{a = 1 \\ \gcd(a,m)=1}}^{m-1}
\overline{\chi}(a) \left\{\frac{na}{m}\right\}
= \chi(n)\, B_{1,\overline{\chi}}.
\]
\end{lemma}

\begin{proof}
The fractional part $\{na/m\}$ is periodic in $a$
modulo~$m$. Since $\gcd(n,m)=1$, the map
$a \mapsto na$ permutes the units, so
\begin{align*}
\sum_a \overline{\chi}(a)\{na/m\}
&= \sum_a \overline{\chi}(n^{-1}a)\{a/m\} \\
&= \chi(n)\sum_a \overline{\chi}(a)\{a/m\}
= \chi(n)\,B_{1,\overline{\chi}},
\end{align*}
where the last step uses
$B_{1,\overline{\chi}}
= (1/m)\sum a\,\overline{\chi}(a)
= \sum \overline{\chi}(a)\,\{a/m\}$
(since $\{a/m\}=a/m$ for $1 \le a \le m{-}1$).
\end{proof}

% ============================================================
\section{The $L$-Function Encoding}

\begin{corollary}\label{cor:L-encoding}
For prime $b$ and primitive odd $\chi$ modulo $b^2$:
\[
|\hat{S}^{\circ}(\chi)|
= \frac{\sqrt{m}}{\pi\, \phi(m)}\,
|L(1, \chi)| \cdot |S_G(\chi)|.
\]
\end{corollary}

\begin{proof}
By the Bernoulli--$L$-value formula for odd primitive
characters~\cite{davenport},
$|B_{1,\overline{\chi}}| = \sqrt{m}\,
|L(1, \chi)| / \pi$.
Substituting into
$|\hat{S}^{\circ}| = |B_1| \cdot |S_G| / \phi$
gives the result.
\end{proof}

The collision coefficient factors into:
\begin{itemize}
\item $|L(1, \chi)|$: the $L$-function special value,
measuring distance from zeros;
\item $|S_G(\chi)|$: the diagonal character sum,
which by the $S_G$ factorization
theorem~\cite{paperC} equals
$2\,|\sum_{k=1}^{b-1} \overline{\chi}(k)|$.
\end{itemize}

% ============================================================
\section{The Double Encoding}

The $S_G$ factorization gives
$|S_G| = 2|\sum_{k=1}^{b-1} \overline{\chi}(k)|$.
Computation shows this partial sum correlates with
$|L(1, \chi)|$.

\begin{table}[h]
\centering
\begin{tabular}{rrrr}
\toprule
$b$ & $n$ & $r(|P|, |L|)$ &
$r \cdot \log b$ \\
\midrule
$5$ & $8$ & $1.00$ & $1.61$ \\
$7$ & $18$ & $0.88$ & $1.71$ \\
$13$ & $72$ & $0.83$ & $2.13$ \\
$19$ & $162$ & $0.86$ & $2.52$ \\
$31$ & $450$ & $0.84$ & $2.90$ \\
$43$ & $882$ & $0.85$ & $3.18$ \\
$61$ & $1800$ & $0.84$ & $3.46$ \\
$71$ & $2450$ & $0.84$ & $3.57$ \\
\bottomrule
\end{tabular}
\caption{Correlation between the partial sum
$|P| = |\sum_{k=1}^{b-1} \overline{\chi}(k)|$ and
$|L(1, \chi)|$ across primitive odd characters.
The correlation is consistent with decay on the
scale of $c / \log b$.}
\label{tab:decay}
\end{table}

At base~$5$, the encoding is exact:
$|P| = (\sqrt{5}/2)\,|B_{1,\overline{\chi}}|$ for
every primitive character, giving
$|\hat{S}^{\circ}| \propto |L(1)|^2$. At larger
bases, the correlation decays, and the effective
exponent in $|\hat{S}^{\circ}| \sim |L(1)|^{\alpha}$
decreases from $2.00$ toward $1$.

% ============================================================
\section{The Apostol Decomposition}

The partial sum $P = \sum_{k=1}^{b-1}
\overline{\chi}(k)$ has an exact decomposition
via the Apostol--Berndt identity for short character
sums~\cite{davenport}. Normalizing by the Gauss sum:
\[
\widetilde{P} = \frac{\pi}{\tau(\chi)}\, P
= L(1, \overline{\chi}) + \Delta(\chi),
\]
where the \emph{packet}
\[
\Delta(\chi) = \frac{i}{\varphi(b)}
\sum_{\substack{\xi \bmod b \\ \xi(-1) = 1,\;
\xi \ne 1}} \tau(\overline{\xi})\,
L(1, \xi\overline{\chi})
\]
is a sum of twisted $L$-values weighted by Gauss
sums over the even characters modulo~$b$.

\begin{observation}[Packet magnitude]
The mean ratio $|\Delta| / |L(1)|$ is constant at
approximately $1.08$ across all prime bases tested
($b = 5$ to $43$). The packet has the same magnitude
as the main term.
\end{observation}

\begin{observation}[Phase uniformity]
The relative phase
$\theta = \arg(\Delta) - \arg(L(1, \overline{\chi}))$
has $\langle \cos\theta \rangle = 0.000$ at every base.
The packet's phase is uniformly distributed relative to
$L(1)$.
\end{observation}

\begin{observation}[Variance decay]
The standard deviation of $|\Delta| / |L(1)|$ across
primitive odd characters is consistent with decay on
the scale of $c / \log b$, with
$\operatorname{std} \cdot \log b \approx 1.2$ across
all bases tested.
\end{observation}

A natural explanation for the correlation decay is
the observed packet behavior. The partial sum $\widetilde{P}
= L(1) + \Delta$ combines two terms of equal
magnitude and uniform relative phase. When the ratio
$|\Delta|/|L|$ has high variance (small $b$), the
variation in $|\widetilde{P}|$ tracks the variation
in $|L|$, producing high correlation. As $b$ grows,
the variance of $|\Delta|/|L|$ decreases (the packet
becomes more uniform relative to $L$), and the
$L$-specific information in $|\widetilde{P}|$ is
washed out.

This is consistent with the fact that $\Delta$ is a sum of
$(\varphi(b)/2 - 1)$ twisted $L$-value terms. As the
number of terms grows, the relative fluctuation of
the sum decreases, concentrating
$|\Delta|/|L|$ around its mean.

% ============================================================
\section{The Moment Theorem}

The decomposition theorem translates the primitive odd
sector of the collision transform into a mean-value
statement about $L$-function special values. Write
\[
F^{\circ}_{\mathrm{prim}}(s)
=
\sum_{\substack{\chi \bmod m \\
\chi \text{ primitive odd}}}
\hat{S}^{\circ}(\chi)\, P(s,\chi),
\qquad m=b^2.
\]

\begin{theorem}[Moment majorant]\label{thm:moment}
For prime $b$ and primitive odd characters modulo
$m=b^2$, one has
\[
\left|F^{\circ}_{\mathrm{prim}}(1)\right|
\le
\frac{1}{\phi(m)}
\sum_{\substack{\chi \bmod m \\
\chi \text{ primitive odd}}}
|B_{1,\overline{\chi}}| \cdot |S_G(\chi)|
\cdot |P(1, \chi)|
\]
where
$P(1, \chi) = \sum_p \chi(p)/p$.
\end{theorem}

\begin{proof}
By the decomposition theorem,
\[
F^{\circ}_{\mathrm{prim}}(1)
= -\frac{1}{\phi}\sum_{\chi\text{ prim.\ odd}}
B_{1,\overline{\chi}}\cdot
\overline{S_G(\chi)}\cdot P(1,\chi).
\]
The triangle inequality gives the stated majorant.
\end{proof}

By the Bernoulli--$L$-value formula,
$|B_1|=\sqrt{m}\,|L(1,\chi)|/\pi$.
By the $S_G$ factorization,
$|S_G|=2\,|\sum_{k=1}^{b-1}\overline{\chi}(k)|$.
And $P(1,\chi)$ relates to $\log L(1,\chi)$ via the
Euler product. The majorant therefore places
$|F^{\circ}_{\mathrm{prim}}(1)|$ beneath an explicit
weighted combination of
$|L(1,\chi)|\cdot|S_G(\chi)|\cdot|\log L(1,\chi)|$
across the primitive odd characters modulo~$m$.

\begin{corollary}
At lag~$1$, the modulus $m = b^2$
has $\phi(m)/2 = b(b{-}1)/2$ odd characters. The
moment majorant expresses
$|F^{\circ}_{\mathrm{prim}}(1)|$ as lying beneath an
explicit weighted average over the primitive odd family.
\end{corollary}

\begin{remark}
At base~$5$, lag~$1$: $|S_G| = (5\sqrt{5}/\pi)
|L(1)|$ exactly, so the moment majorant reads
\[
\left|F^{\circ}_{\mathrm{prim}}(1)\right|
\le C \sum |L(1)|^2 |\log L(1)|
\]
for an explicit constant $C$: a weighted second
moment majorizing the finite collision transform at
the edge.
\end{remark}

\begin{theorem}[Conditional cross-moment]
\label{thm:cross}
Fix $\sigma_0 \in (1/2, 1)$ and suppose every primitive
odd $L$-function modulo $m=b^2$ satisfies
$L(s, \chi) \ne 0$ for
$\operatorname{Re}(s) > \sigma_0$.
Then for every real $s > \sigma_0$:
\[
\left|F^{\circ}_{\mathrm{prim}}(s)\right|
\le
\frac{1}{\phi(m)}
\sum_{\substack{\chi \bmod m \\
\chi \text{ primitive odd}}}
|B_{1,\overline{\chi}}| \cdot |S_G(\chi)|
\cdot |P(s, \chi)|.
\]
\end{theorem}

\begin{proof}
By the conditional penetration
theorem~\cite{paperB}, $F^{\circ}(s)$ converges for
real $s > \sigma_0$. The decomposition theorem gives
$F^{\circ}_{\mathrm{prim}}(s) = -(1/\phi) \sum B_1 \cdot
\overline{S_G} \cdot P(s, \chi)$. The triangle
inequality gives the bound.
\end{proof}

Since $|B_1| = \sqrt{m}\, |L(1, \chi)| / \pi$ and
$P(s, \chi) = \log L(s, \chi) - H(s, \chi)$, the
cross-moment majorant places $|F^{\circ}_{\mathrm{prim}}(s)|$
beneath a weighted sum of products
\[
|L(1, \chi)| \cdot |S_G(\chi)| \cdot
|\log L(s, \chi)|
\]
across odd characters. This connects $L$-values at
two points: $s = 1$ (the edge, through $B_1$) and
$s = \sigma_0$ (in the strip, through $P$). As
$\sigma_0 \to 1/2$, the bound constrains
$L$-values approaching the critical line, weighted
by $L$-values at the edge and filtered by the
collision geometry through $S_G$.

% ============================================================
\section{The Zero Coupling}

The decomposition theorem gives the exact structural
identity behind part of what the preceding
papers~\cite{paper19,paper20} observed. The factor
$B_1$ shows that characters with small $|L(1,\chi)|$
necessarily have small collision coefficients
$|\hat{S}^{\circ}(\chi)|$. This does not by itself
prove the anti-correlation, the spectral repulsion,
or the double transversality, since those involve
additional information about $S_G(\chi)$, partial
prime sums, and cross-character geometry. It does,
however, explain why the edge value $L(1,\chi)$ is a
structural input to every coefficient.

The analytic continuation
from~\cite{paperB} gives
\[
\mathcal{F}^{\circ}(s)
= -\frac{1}{\phi} \sum_{\chi \text{ odd}}
B_{1,\overline{\chi}} \cdot \overline{S_G(\chi)}
\cdot \bigl[\log L(s, \chi) - H(s, \chi)\bigr].
\]
Near a zero $\rho$ of $L(s, \chi_0)$, the term
$\log L(s, \chi_0) \to -\infty$. The coefficient
of this divergence is
$B_{1,\overline{\chi}_0} \cdot
\overline{S_G(\chi_0)} / \phi$, which involves
$L(1, \chi_0)$ through $B_1$. The behavior of the
collision transform near a zero at $s = \rho$ is
therefore controlled by the $L$-value at $s = 1$.

\begin{remark}
The digit function couples different points of the
critical strip: the singularity structure at depth
$s$ is weighted by the health at $s = 1$. A zero of
$L(s, \chi)$ deep in the strip produces a
singularity in $\mathcal{F}^{\circ}$ whose residue
is proportional to $|L(1, \chi)|$. Characters that
are healthy at the edge produce large residues at
depth. Characters that are weak at the edge produce
small residues. The collision geometry mediates this
coupling through the Bernoulli number.
\end{remark}

% ============================================================
\section{Remarks}

The decomposition
$\hat{S}^{\circ} = -B_1 \cdot \overline{S_G} / \phi$
is exact for every primitive odd character at every
prime base. It connects the collision invariant to
two classical objects: the Bernoulli number $B_1$
(encoding $|L(1)|$ via the Bernoulli--$L$-value formula) and
the partial character sum $P = \sum_{k=1}^{b-1}
\overline{\chi}(k)$ (encoding a second, decaying
factor of $|L|$).

At base~$5$, the double encoding is exact:
$|\hat{S}^{\circ}| \propto |L(1)|^2$. At larger
bases, the second factor's encoding degrades
logarithmically. The digit function's ``opening
act'' (its values at the first $b{-}1$ units) carries
information about the $L$-function landscape, but the
fidelity of this information is consistent with
decay on the scale of $1/\log b$ across the tested
prime bases.

For imaginary quadratic fields $\mathbb{Q}(\sqrt{-d})$,
the class number satisfies
$h(-d) = (\sqrt{d}/\pi)\, |L(1, \chi_d)|$.
The decomposition theorem therefore encodes class
numbers: at base~$5$,
$|\hat{S}^{\circ}| \propto h(-d)^2$. The digit
function squares class numbers.

The open problem is to prove this decay rate. The
decomposition theorem, the Abel summation formula,
and the near-linear structure of $G(k)$ provide the
algebraic framework. What remains is a bound on the
correlation between partial character sums over
$\{1, \ldots, b{-}1\}$ and $L$-function special
values for primitive characters modulo~$b^2$.

% ============================================================
\begin{thebibliography}{99}

\bibitem{davenport}
H.~Davenport,
\emph{Multiplicative Number Theory},
3rd~ed., Springer, 2000.

\bibitem{iwaniec}
H.~Iwaniec and E.~Kowalski,
\emph{Analytic Number Theory},
AMS Colloquium Publications, vol.~53, 2004.

\bibitem{paperA}
A.~S.~Petty,
The collision invariant,
preprint, 2026.

\bibitem{paperB}
A.~S.~Petty,
The collision transform,
preprint, 2026.

\bibitem{paperC}
A.~S.~Petty,
The collision spectrum,
preprint, 2026.

\bibitem{paper19}
A.~S.~Petty,
The general neutrality theorem,
research note, 2025.

\bibitem{paper20}
A.~S.~Petty,
The double transversality,
research note, 2025.

\bibitem{nfield}
A.~S.~Petty,
\emph{nfield: A structural analysis engine for fractional
field invariants},
software.
\url{https://github.com/alexspetty/nfield}

\end{thebibliography}

\end{document}
